% =============================================================================
% Geometric Competition as a Shared Principle of Ambiguity
% in LLM and Neural Emotion Representations
%
% NeurIPS 2026 Submission — Double-blind review version
%
% DEPENDENCIES:
%   neurips_2024.sty  — download from https://neurips.cc/Conferences/2024/PaperInformation/StyleFiles
%                       and place in the same directory as this .tex file.
%   references.bib    — provided alongside this file.
%
% COMPILE:
%   pdflatex neurips_submission.tex
%   bibtex neurips_submission
%   pdflatex neurips_submission.tex
%   pdflatex neurips_submission.tex
% =============================================================================

\documentclass{article}

\usepackage[preprint]{neurips_2024}   % change to [final] for camera-ready

\usepackage[utf8]{inputenc}
\usepackage[T1]{fontenc}
\usepackage{microtype}
\usepackage{amsmath,amssymb}
\usepackage{graphicx}
\usepackage{subcaption}
\usepackage{booktabs}
\usepackage{xcolor}
\usepackage{hyperref}
\usepackage{url}
\usepackage{multirow}

% ---------------------------------------------------------------------------
% Figure search paths (relative to the location of this .tex file)
% ---------------------------------------------------------------------------
\graphicspath{
  {../../experiments/understanding_text_embeddings/reports/phase1/}
  {../../experiments/understanding_text_embeddings/reports/phase2/}
  {../../experiments/understanding_text_embeddings/reports/phase3/}
  {../../experiments/understanding_text_embeddings/reports/phase4/}
  {../../experiments/understanding_text_embeddings/reports/phase5/}
  {../../experiments/brain_embedding_understanding/global_behavior_comparison/}
  {../../experiments/brain_embedding_understanding/checking_context_retention_across_dimensions/reports/}
  {../../experiments/brain_embedding_understanding/checking_centroids_with_spatial_context_data/reports/}
  {../../experiments/brain_embedding_understanding/checking_centroids/reports/}
  {../../experiments/brain_embedding_understanding/valence-arousal-dimensional_reduction/reports/}
  {../../experiments/brain_embedding_understanding/checking_density_geometry/reports/}
  {../../experiments/understanding_text_embeddings/reports/phase4/visuals/}
}

% ---------------------------------------------------------------------------
\title{Beyond Distance: Geometric Competition Explains Ambiguity in Neural and LLM Representations}

\author{
  Pritish Ranjan Verma, Daniel Sikar, Aimee Bottrill-Frost, Joshua Narindra Bhawanlall \\
  \textit{Submitted to NeurIPS 2026}
}

\begin{document}
\maketitle

% ===========================================================================
\begin{abstract}
% ===========================================================================
A widely used geometric heuristic is that uncertainty is governed by radial distance to
class prototypes: points near centroids are treated as certain, while distant points are
treated as ambiguous.
We show this account is incomplete.
Studying emotion representations across eight large language model (LLM) variants and
fMRI from 40 human subjects, we find that both artificial and biological systems organise
emotional content in a consistent radial structure --- a dense shell of observed instances
at intermediate distances from class centroids, with reduced density in the immediate
vicinity of the centroids.
Critically, ambiguity within this structure is not determined by raw distance to a prototype
but by relative competition between prototypes.
We formalise this as centroid margin and demonstrate that it predicts classifier uncertainty
in brain representations substantially better than raw distance, with the relationship
collapsing under feature shuffling.
In fine-tuned LLMs, geometric margin is tightly coupled to logit confidence.
We further find that emotional content is encoded in a low-rank linear subspace occupying
a small fraction of embedding dimensions, with additional dimensions contributing limited
signal.
Despite this shared local geometry, global representational organisation is inverted:
LLMs are structured primarily along valence while brain representations are structured
primarily along arousal, producing near-opposite dissimilarity matrices that strengthen
as biological noise is reduced.
These findings establish a reusable geometric framework for comparing representational
ambiguity across artificial and neural systems, while revealing a fundamental divergence
in how the two systems prioritise the affective dimensions of emotional space.
\end{abstract}

% ===========================================================================
\section{Introduction}
% ===========================================================================

How emotion is represented geometrically in artificial and biological neural systems is
a foundational question at the intersection of representation learning and affective
neuroscience. While large language models have achieved strong performance on emotion
classification, the internal geometry supporting these classifications remains poorly
characterised. The fMRI literature has documented that emotions are encoded in
distributed cortical patterns~\citep{horikawa2020}, but the geometric relationship
between those patterns and the structured embedding spaces of language models has not
been systematically investigated.

This paper addresses a specific and tractable version of this question: do LLM embedding
spaces and human fMRI activation spaces share a geometric principle for representational
ambiguity? We operationalise ambiguity through \emph{centroid distance} and
\emph{centroid margin} --- the competitive gap between a point's distance to its own class
prototype versus the nearest competing prototype. Our hypothesis is that both systems
exhibit a structured void near class prototypes, a peak-density belt of real instances
at moderate distances, and a gradient from certainty to ambiguity as geometric exposure
to competing prototypes increases.

\textbf{Research gap.}
Prior work on LLM--brain alignment has focused on language-processing tasks using RSA
or encoding models~\citep{caucheteux2022,schrimpf2021,li2023b,huth2016,li2023a},
relating brain regions to model layers.
Emotion representation geometry has been studied in affective
computing~\citep{russell1980,koelstra2012} but rarely with the aim of characterising
the full distributional structure of embedding spaces across a void--belt--margin
framework. Whether the ambiguity gradient is a shared principle across artificial and
biological systems has not been previously investigated.

\textbf{Contributions.}
\begin{enumerate}
  \item A normalised-distance framework enabling geometric comparison across spaces of
        different dimensionality, identifying the Void and the Belt as consistent
        structural features across all tested systems.
  \item Characterisation of overlap rates and their pairwise structure across eight
        LLM variants.
  \item \textbf{The Low-Rank Emotion Subspace}: emotion is a low-rank linear property
        occupying $\approx\!2.6\%$ of embedding capacity (20 of 768 dimensions).
        Accuracy in 20D equals or exceeds 768D across all tested architectures.
  \item A margin-vs-distance diagnostic distinguishing relational from radial geometric
        structure: centroid \emph{margin} predicts brain uncertainty
        (AUC$\,{=}\,0.81$); centroid \emph{distance} fails (AUC$\,{=}\,0.56$).
  \item The global representational inversion: LLMs sort emotions by valence; the brain
        by arousal. This produces near-opposite dissimilarity matrices that deepen as
        biological noise is reduced.
\end{enumerate}

% ===========================================================================
\section{Related Work}
% ===========================================================================

\textbf{Representational geometry in LLMs.}
\citet{mikolov2013} showed that distributional embeddings encode semantic structure in
linear subspaces.
\citet{ethayarajh2019} documented anisotropy in contextualised representations.
\citet{kumar2022} showed fine-tuning can compress pretrained features --- directly
relevant to our compaction finding.
Our contribution extends this work by characterising the full radial density and
competitive margin geometry of emotion embeddings.

\textbf{Brain--LLM alignment.}
\citet{caucheteux2022} showed partial convergence between brain activations and LLMs
during language processing.
\citet{schrimpf2021} characterised neural architecture convergence with transformer
predictive processing.
\citet{huth2016} linked distributional semantics to cortical topography.
\citet{ismayilzada2026} found alignment during creative thinking.
Our work differs: we study emotion rather than language processing, focus on the local
geometric ambiguity mechanism, and document a systematic global inversion.

\textbf{Distributed neural coding of emotion.}
\citet{horikawa2020} showed visually evoked emotion representations are high-dimensional,
categorical, and distributed across transmodal brain regions.
\citet{haxby2001} established that cognitive categories are encoded as distributed
activation patterns --- the biological prior explaining why brain fMRI shows negative
silhouette in 2D projection while remaining decodable in 48 dimensions.

\textbf{The circumplex model.}
\citet{russell1980} organised affective states along valence and arousal dimensions,
operationalised through DEAP~\citep{koelstra2012}, ANEW~\citep{bradley1999}, and
IAPS~\citep{lang2008}.
We use external V-A reference coordinates from this literature to interpret global
representational axes, not as measured participant ground truth.

% ===========================================================================
\section{Datasets and Models}
% ===========================================================================

\textbf{Text emotion dataset.}
We use the \texttt{dair-ai/emotion} dataset~\citep{saravia2018}, a 6-class text emotion
corpus (anger, fear, happiness, love, sadness, surprise).
We use the validation split, balanced to \textbf{4,038 samples} (673 per class).
Class balance ensures centroids are not displaced by class-size differences and that
overlap rates are comparable across categories.
Labels are integer-encoded: anger=0, fear=1, happiness=2, love=3, sadness=4, surprise=5.

\textbf{Brain fMRI dataset.}
We use OpenNeuro DS005700~\citep{openneuro2024}, ``Neural MO --- fMRI Dataset for
Emotion Recognition.''
The dataset contains \textbf{40 subjects} viewing emotion-eliciting stimuli under five
conditions: afraid, calm, delighted, depressed, and excited.
Each subject contributes activations across \textbf{48 cortical ROIs} from the
Harvard-Oxford atlas, yielding 200 observations total (40 subjects $\times$ 5 emotions).
Repeated blocks per subject and emotion are averaged to stable ROI vectors.
Brain LOSO decoding accuracy is 0.56 (95\% CI: 0.49--0.63, permutation $p<0.001$),
well above the 5-class chance level of 0.20.
Our $N\!=\!40$ exceeds high-impact precedents with smaller samples:
Individual Brain Charting ($N\!=\!12$)~\citep{openneuro2024},
Precision Functional Mapping of Children ($N\!=\!12$), and
Natural Scenes Dataset ($N\!=\!8$).

\textbf{LLM variants.}
We evaluate eight variants by crossing two base models, two training states, and two
extraction layers:

\begin{center}
\small
\begin{tabular}{lllll}
\toprule
Variant & Base Model & State & Layer & Dim \\
\midrule
MPNet-FT-Final  & all-mpnet-base-v2~\citep{song2020} & Fine-tuned & L12 & 768 \\
MPNet-FT-Mid    & all-mpnet-base-v2 & Fine-tuned & L6  & 768 \\
MPNet-Base-Final& all-mpnet-base-v2 & Pretrained & L12 & 768 \\
MPNet-Base-Mid  & all-mpnet-base-v2 & Pretrained & L6  & 768 \\
BGE-FT-Final    & bge-base-en-v1.5~\citep{xiao2023} & Fine-tuned & L12 & 768 \\
BGE-FT-Mid      & bge-base-en-v1.5 & Fine-tuned & L6  & 768 \\
BGE-Base-Final  & bge-base-en-v1.5 & Pretrained & L12 & 768 \\
BGE-Base-Mid    & bge-base-en-v1.5 & Pretrained & L6  & 768 \\
\bottomrule
\end{tabular}
\end{center}

Fine-tuned models were trained on the \texttt{dair-ai/emotion} training split.
The cross-system context retention experiment (Section~\ref{sec:finding2}) additionally
includes Qwen3-1.7B fine-tuned~\citep{qwen2025} (PCA-projected to 768D).

% ===========================================================================
\section{Methods}
% ===========================================================================

\textbf{Embedding extraction.}
All text embeddings use attention-mask-weighted mean pooling:
\begin{equation}
  \mathbf{e} = \frac{\sum_{i=1}^{N} \mathbf{h}_i \cdot m_i}{\sum_{i=1}^{N} m_i}
  \label{eq:pooling}
\end{equation}
where $\mathbf{h}_i$ is the hidden state for token $i$ and $m_i$ is the attention mask.
Final-layer (L12) extraction uses the last hidden state directly.
Middle-layer (L6) extraction enables \texttt{output\_hidden\_states=True} and accesses
\texttt{hidden\_states[6]}.
Tokenisation: model-specific tokenisers with \texttt{max\_length=128},
\texttt{padding=True}, \texttt{truncation=True}, batch size 32.

\textbf{Normalised distance framework.}
To compare across systems with different intrinsic scales (768D LLMs vs.\ 48D brain
ROI), all distances are normalised by the mean radius of each system:
\begin{equation}
  \tilde{d}(x, c_k) = \frac{d(x, c_k)}{\bar{R}_k},
  \qquad
  \bar{R}_k = \frac{1}{N_k}\sum_{x \in C_k} d(x, c_k)
  \label{eq:normdist}
\end{equation}
This maps all systems to a common scale ($\approx\!0$--$2.5$ normalised units) without
discarding distributional shape.
A point geometrically overlaps class $C'$ (its non-assigned class) if
$d(P, c_{C'}) < d(P, c_C)$.
The \emph{centroid margin} is:
\begin{equation}
  m(P) = d(P,\, c_{\text{nearest competitor}}) - d(P,\, c_C)
  \label{eq:margin}
\end{equation}
Positive margin: $P$ is closer to its true centroid than to any competitor (more
certain). Negative margin: geometric overlap (more ambiguous).

\textbf{Brain 48D$\to$11D transformation.}
To test whether the global representational inversion (Section~\ref{sec:inversion})
is noise-driven, we construct a domain-informed 11-dimensional representation by
aggregating ROI activations into biologically coherent groups.
The Default Mode Network (DMN)~\citep{li2014dmn} supports self-referential processing
and social cognition.
The Salience Network~\citep{seeley2019} mediates detection of affectively significant
stimuli.
The Central Executive Network~\citep{miller2001} coordinates regulatory responses to
emotional content.
The 11 dimensions consist of five anatomical lobe means (frontal, temporal, parietal,
occipital, cingulate/limbic), five functional network means (DMN, Salience, Central
Executive, visual, somatomotor), and one neighbour-context feature capturing local
inter-regional coordination.
This is domain-informed structured denoising, not generic PCA: each dimension
corresponds to a biologically meaningful functional unit.
Classification performance in 11D remains above chance and comparable to 48D,
confirming the transformation preserves task-relevant signal.

\textbf{Valence-arousal reference mapping.}
External V-A coordinates are drawn from the circumplex model~\citep{russell1980,
mehrabian1974} and affective-norm resources DEAP~\citep{koelstra2012},
ANEW~\citep{bradley1999}, and IAPS~\citep{lang2008} (1--9 SAM scale).
Because the fMRI dataset supplies categorical labels rather than participant-level V-A
ratings, these coordinates function as an \emph{external reference geometry},
not measured subjective ground truth.
Sensitivity analysis (Appendix~\ref{app:va-sensitivity}) perturbs coordinates with
Gaussian noise (SD\,=\,0.25/0.50/0.75) to verify qualitative robustness.

\textbf{RSA methodology.}
We follow \citet{kriegeskorte2008}.
For each system, pairwise Euclidean distances between emotion class centroids form a
$K \times K$ Representational Dissimilarity Matrix (RDM).
Brain--LLM comparison uses the 3-emotion overlap: fear$\approx$afraid,
happiness$\approx$delighted, sadness$\approx$depressed.
Systems are compared by Pearson correlation between vectorised RDMs.
A brain--brain noise ceiling estimates the upper bound imposed by individual subject
variability: upper 0.484, lower 0.460 (mean $\approx\!0.47$), computed by split-half
resampling of the 40-subject dataset.
With $K\!=\!3$ emotions, RDMs contain 3 unique pairwise distances;
reported Pearson correlations have 1 degree of freedom.
Specific magnitudes are therefore indicative; directional findings (sign, contrast with
noise ceiling) are the reliable quantities.

% ===========================================================================
\vspace{-1.5em}
\section{Results}
\vspace{-0.5em}
% ===========================================================================

We present five findings in order of increasing cross-system scope.
Findings 1--3 characterise LLM geometry; Finding 4 introduces the ambiguity gradient;
Finding 5 documents the global representational inversion.

% ---------------------------------------------------------------------------
\vspace{-0.5em}
\subsection{Finding 1: The Void and the Belt}
\label{sec:finding1}
\vspace{-0.3em}
% ---------------------------------------------------------------------------

A consistent two-zone density structure emerges across all 8 LLM variants and
human brain fMRI (Figure~\ref{fig:clusters-density}).
The structural parallel is qualitative: packing density magnitudes differ substantially
across systems, and density-shape similarity does not reach conventional significance
under permutation testing (Wasserstein $p\!=\!0.46$).
The claim is descriptive --- shared organisational logic, not matched compression.

\textbf{The Void.}
No data points exist within approximately 0.375--0.625 normalised distance units of
any emotion centroid.
Fine-tuned models show first non-zero density at $\approx\!0.4375$ units;
pretrained models at $\approx\!0.6875$ units.
The Void arises as a geometric consequence of distributed within-class variance:
a centroid is the mean of a distributed ring of instances, so it sits inside the ring
in zero-density space.

\textbf{The Belt.}
Density rises sharply from the void boundary and peaks at approximately 0.56--0.94
normalised units.
Fine-tuned models peak earlier (MPNet-FT-Final at 0.5625, BGE-FT-Final at 0.6875);
pretrained models peak later (both at 0.9375).

\vspace{0.4em}
\noindent
\begin{minipage}[t]{0.38\textwidth}
  \centering
  \includegraphics[height=4.5cm]{global_density_ambiguity_curves}
  \par\vspace{0.2em}
  {\footnotesize (a) Radial density profiles (top) and ambiguity gradients
  (bottom) for all 8 LLM variants. The void region (zero density near
  centroid) and the Belt (peak density shell) are consistent across all systems.}
\end{minipage}%
\hfill%
\begin{minipage}[t]{0.59\textwidth}
  \centering
  \includegraphics[height=4.5cm]{global_density_decay_comparison_}
  \par\vspace{0.2em}
  {\footnotesize (b) Radial density decay for LLM variants (coloured)
  overlaid with human brain fMRI (dashed black).
  The two-zone structure extends to neural data.}
\end{minipage}
\refstepcounter{figure}\label{fig:clusters-density}
\par\vspace{0.3em}
\noindent{\small\textbf{Figure~\thefigure: The Void and the Belt (Finding 1).}
(a) All 8 LLM variants share the same two-zone radial density structure
in normalised-distance space.
(b) The same structure is present in human brain fMRI emotion representations,
establishing it as a cross-system geometric principle.}
\par\vspace{0.5em}

\textbf{Certainty buffer.}
The radial gap between density peak and onset of geometric ambiguity
($>\!5\%$ overlap rate; geometric ambiguity defined as the region where
nearest-neighbour class boundaries overlap in embedding space):

\begin{center}
\small
\begin{tabular}{lcccc}
\toprule
Variant & Void Onset & Density Peak & Overlap Onset & Certainty Buffer \\
\midrule
MPNet-FT-Final  & 0.4375 & 0.562 & 2.438 & $+1.875$ \\
BGE-FT-Final    & 0.4375 & 0.688 & 2.438 & $+1.750$ \\
MPNet-Base-Final& 0.8125 & 0.938 & 0.938 & $\phantom{+}0.000$ \\
BGE-Base-Final  & 0.6875 & 0.938 & 1.062 & $+0.125$ \\
\bottomrule
\end{tabular}
\end{center}

Fine-tuning creates a geometrically safe zone between peak density and the onset of
competitive class confusion; for pretrained models this buffer is near zero.
No embedding aligns exclusively with a single class prototype in any system:
global overlap rates range from 1.76\% (MPNet-FT-Final) to 19.22\% (BGE-Base-Final),
confirming that fine-tuning substantially reduces inter-class geometric confusion.

% ---------------------------------------------------------------------------
\subsection{Finding 2: Information Compaction vs.\ Distribution}
\label{sec:finding2}
% ---------------------------------------------------------------------------

Iterative SVD ablation removes the most informative linear direction at each step;
a fresh logistic regression is re-trained on the residual.
\textbf{Signal erasure points} (dimension at which accuracy reaches the 6-class chance
baseline of 16.7\%):

\begin{center}
\small
\begin{tabular}{lcc}
\toprule
Variant & Baseline Accuracy (\%) & Erasure Point \\
\midrule
MPNet-FT-Final   & 98.24 & dim 15 \\
BGE-FT-Final     & 97.99 & dim 20 \\
BGE-Base-Final   & 93.78 & dim 26 \\
MPNet-Base-Final & 94.08 & dim 67 \\
\bottomrule
\end{tabular}
\end{center}

Fine-tuned models achieve higher baseline accuracy but concentrate all emotional signal
into 15--20 directions; remove them and signal collapses entirely.
Pretrained MPNet-Base retains decodable signal for 67 dimensions, a $4.5\times$ wider
manifold than the fine-tuned equivalent (Figure~\ref{fig:signal-cliff}).

\begin{figure}[h!]
  \centering
  \begin{subfigure}[t]{0.50\textwidth}
    \includegraphics[height=4cm]{early_decay_comparison_top25}
    \caption{Signal cliff: accuracy as dominant directions are removed iteratively.
             Fine-tuned models reach chance at dimensions 15--20;
             pretrained BGE at 26, pretrained MPNet at 67.}
    \label{fig:signal-cliff}
  \end{subfigure}
  \hfill
  \begin{subfigure}[t]{0.48\textwidth}
    \includegraphics[height=4cm]{Content_Retention_Comparison_}
    \caption{Cross-system context retention.
             Brain $D_{50}\!=\!4$; fine-tuned LLM $D_{50}\!=\!12$;
             pretrained LLM $D_{50}\!=\!17$.}
    \label{fig:context-retention}
  \end{subfigure}
  \caption{\textbf{Signal structure across systems (Finding 2).}
           (a) SVD ablation reveals qualitatively different compaction strategies.
           (b) The brain requires only 4 dominant directions to achieve half its
           classification accuracy, yet maintains a distributed plateau rather than
           collapsing to chance --- consistent with distributed biological coding.}
  \label{fig:signal-structure}
\end{figure}

\textbf{Cross-system extension.}
A separate iterative ablation across three systems
(Figure~\ref{fig:context-retention}) shows qualitatively distinct signal geometries:
the human brain ($D_{50}\!=\!4$, AUC 0.302) uses the most distributed code ---
removing just 4 dominant directions halves accuracy, yet the curve flattens into a
plateau rather than collapsing to chance~\citep{haxby2001};
Qwen-FT ($D_{50}\!=\!12$, AUC 0.260) compresses signal into a steep
cliff; pretrained MPNet ($D_{50}\!=\!17$, AUC 0.331) spreads it across the widest
manifold.

% ---------------------------------------------------------------------------
\subsection{Finding 3: The Low-Rank Emotion Subspace}
\label{sec:finding3}
% ---------------------------------------------------------------------------

\textbf{Named result.}
Emotion is a low-rank linear property, occupying $\approx\!20/768\,{\approx}\,2.6\%$
of full embedding capacity.
The cloud appearance of pretrained models in 768D is not an absence of geometry ---
it is the geometry of $\approx\!748$ noise dimensions overwhelming 20 signal
dimensions.

\textbf{Silhouette scores and classification accuracy before and after projection into
the top 20 SVD directions (Figure~\ref{fig:silhouette}):}

\begin{center}
\small
\setlength{\tabcolsep}{4pt}
\begin{tabular}{lcccccc}
\toprule
Variant & 768D Sil. & 20D Sil. & Sil.\ Change & 768D Acc.\ (\%) & 20D Acc.\ (\%) \\
\midrule
MPNet-FT-Final   & 0.860 & 0.907 & $+5.5\%$  & 98.24 & 98.27 \\
BGE-FT-Final     & 0.685 & 0.835 & $+22.0\%$ & 97.99 & 98.09 \\
MPNet-Base-Final & 0.047 & 0.408 & $+765\%$  & 94.08 & 95.64 \\
BGE-Base-Final   & 0.057 & 0.415 & $+627\%$  & 93.78 & 95.12 \\
\bottomrule
\end{tabular}
\end{center}

The extra 748 dimensions contribute only noise that slightly suppresses linear
classification performance.
Pairwise centroid distances preserve the same relational structure in 20D as in 768D:
Happiness--Love and Fear--Anger are the closest pairs;
Sadness--Surprise and Happiness--Fear are the most distant.

\enlargethispage{6\baselineskip}
\begin{figure}[h!]
  \centering
  \includegraphics[width=0.60\textwidth]{silhouette_comparison_20D}
  \caption{\textbf{The Low-Rank Emotion Subspace (Finding 3).}
           Silhouette scores in full 768D (grey) vs.\ the isolated top-20 SVD subspace
           (red). Pretrained models improve from near-zero to $\approx\!0.41$;
           fine-tuned models improve modestly from already-high baselines.
           Accuracy in 20D equals or exceeds 768D in all four variants (see text).}
  \label{fig:silhouette}
\end{figure}

\begin{figure}[h!]
  \centering
  \begin{subfigure}[t]{0.49\textwidth}
    \includegraphics[height=3cm]{comparison_scatter_BGE-Base-Final}
    \caption{BGE-Base-Final: PCA scatter in full 768D (left) vs.\ top-20 SVD
             subspace (right). Diffuse clouds sharpen into separated clusters.}
    \label{fig:scatter-bge}
  \end{subfigure}
  \hfill
  \begin{subfigure}[t]{0.49\textwidth}
    \includegraphics[height=3cm]{comparison_scatter_MPNet-Base-Final}
    \caption{MPNet-Base-Final: same comparison. Silhouette improves from
             0.057 to 0.415 ($+627\%$) and 0.047 to 0.408 ($+765\%$).}
    \label{fig:scatter-mpnet}
  \end{subfigure}
  \caption{\textbf{768D vs.\ top-20 SVD subspace (Finding 3).}
           Projecting into the 20-dimensional signal subspace reveals the cluster
           structure hidden by noise dimensions in the full embedding space.}
  \label{fig:scatter-comparison}
\end{figure}

% ---------------------------------------------------------------------------
\clearpage
\subsection{Finding 4: The Cross-System Ambiguity Gradient}
\label{sec:finding4}
% ---------------------------------------------------------------------------

Both systems show that geometric positioning relative to class prototypes predicts
classification confidence --- through mechanistically distinct constructs.

\textbf{LLMs: centroid margin $\to$ logit margin.}
For fine-tuned LLMs, centroid margin is a strong predictor of logit margin
(Figure~\ref{fig:dist-logit}):

\begin{center}
\small
\begin{tabular}{lcc}
\toprule
Variant & Margin-logit Pearson $r$ & Logit agreement in overlap (\%) \\
\midrule
BGE-FT-Final    & 0.9572 & 93.4  \\
MPNet-FT-Final  & 0.9884 & 100.0 \\
\bottomrule
\end{tabular}
\end{center}

The near-linear relationship ($r\!=\!0.9884$ for MPNet-FT) means centroid margin
almost perfectly determines logit margin.
In overlap regions, logit values are biased toward the competing class 93.4--100\% of
the time.
Note that in fine-tuned models this relationship is partially expected by the training
objective; it is reported as a mechanistic characterisation.

\textbf{Brain: centroid margin $\to$ classifier uncertainty.}
For brain fMRI, raw centroid distance is a weak predictor of classification uncertainty
(Spearman $r\!=\!0.1509$, AUC$\,{=}\,0.5627$ --- near chance).
The correct predictor is centroid margin (Equation~\ref{eq:margin}):
\begin{equation*}
  \text{AUC}_{\text{margin}} = 0.8124
  \qquad\text{vs.}\qquad
  \text{AUC}_{\text{raw distance}} = 0.5627
\end{equation*}
with Spearman $r\!=\!0.6108$, $p\!=\!7.78\!\times\!10^{-22}$.
A feature-shuffle control ($r\!=\!0.0287$, $p\!=\!0.687$) confirms the relationship
vanishes under random permutation.

\begin{figure}[h!]
  \centering
  \begin{subfigure}[t]{0.54\textwidth}
    \includegraphics[height=4.5cm]{dist_logit_scatter_MPNet-FT-Final}
    \caption{LLM: centroid margin vs.\ logit margin for MPNet-FT-Final
             ($r\!=\!0.988$). Geometric margin almost perfectly determines
             logit margin.}
    \label{fig:dist-logit}
  \end{subfigure}
  \hfill
  \begin{subfigure}[t]{0.44\textwidth}
    \includegraphics[height=4.5cm]{Uncertainity_report_brain_data}
    \caption{Brain fMRI: centroid margin predicts classifier uncertainty
             (AUC\,=\,0.8124); raw centroid distance fails
             (AUC\,=\,0.5627). Margin is the operative geometric construct.}
    \label{fig:brain-uncertainty}
  \end{subfigure}
  \caption{\textbf{Geometric predictors of confidence (Finding 4).}
           (a) In LLMs, centroid margin tightly predicts logit margin.
           (b) In the brain, centroid \emph{margin} --- not distance --- predicts
           uncertainty, revealing competition between prototypes as the key signal.}
  \label{fig:finding4}
\end{figure}

\textbf{The ambiguity gradient comparison.}
The analysis interpolates brain and LLM overlap-vs-distance curves onto a shared
normalised-radius grid of 100 equally spaced points ($\tilde{d}\in[0, 2.5]$) and
computes Pearson correlation between the resulting 100-point vectors.
The ambiguity-gradient curves show high shape correspondence
($r\!=\!0.9565$, 95\% CI [0.9370, 0.9725]),
supported by a full-pipeline permutation test ($n\!=\!5000$ shuffles, $p<0.001$)
that destroys cross-system alignment when labels are shuffled.
The $r$ is computed over $N\!=\!100$ interpolated grid points and should not be
treated as an independent-sample $p$-value; the permutation test is the inferential
basis.
KS statistic $= 0.57$ and Cohen's $d\!=\!1.34$ confirm that while gradient
\emph{shape} is shared, absolute packing density differs substantially between systems
--- shared ambiguity structure, distinct representational compaction.

% ---------------------------------------------------------------------------
\subsection{Finding 5: The Valence-Arousal Inversion}
\label{sec:inversion}
% ---------------------------------------------------------------------------

The most novel finding of the study is that the shared local ambiguity principle
coexists with a global inversion of representational geometry.

\textbf{(5a) RDM comparison in 48D raw ROI space (Figure~\ref{fig:rdms}).}

\begin{figure}[h!]
  \centering
  \begin{subfigure}[b]{0.48\textwidth}
    \includegraphics[width=\textwidth]{rdm_cosine_brain}
    \caption{Brain RDM (cosine): arousal-structured. Depressed--calm (0.72) and
             afraid--excited (0.82) are close; affect separates by activation
             intensity.}
    \label{fig:rdm-brain}
  \end{subfigure}
  \hfill
  \begin{subfigure}[b]{0.48\textwidth}
    \includegraphics[width=\textwidth]{rdm_cosine_MPNet-Balanced}
    \caption{LLM RDM (cosine): valence-structured. Sadness--happiness (0.36),
             love--happiness (0.26) are close; affect separates by hedonic polarity.}
    \label{fig:rdm-llm}
  \end{subfigure}
  \caption{\textbf{RDM comparison (Finding 5a).}
           Brain and LLM emotion geometry is near-opposite: the brain organises
           affect by arousal intensity while the LLM organises by hedonic polarity
           (cosine RDM Pearson $r\!=\!-0.8756$; noise ceiling $\approx\!+0.47$).}
  \label{fig:rdms}
\end{figure}

\begin{center}
\small
\begin{tabular}{lcc}
\toprule
Comparison & Pearson $r$ & Notes \\
\midrule
Brain vs.\ fine-tuned LLM (cosine)   & $-0.8756$ & Near-opposite geometry \\
Brain vs.\ pretrained LLM (cosine)   & $-0.5476$ & Confirmed \\
Brain vs.\ fine-tuned LLM (Manhattan)& $-0.5476$ & Metric-invariant validation \\
Brain-brain noise ceiling             & $+0.47$   & Upper 0.484, lower 0.460 \\
\bottomrule
\end{tabular}
\end{center}

RDM comparisons provide directional evidence of global inversion, but because only
three shared emotion categories are available, exact correlation magnitudes are
unstable (bootstrap 95\% CI: $[-0.9967,\,+0.6994]$, as expected with $K\!=\!3$, 1 df).
The noise ceiling contextualises the directionality: brain subjects agree with each
other at $r\!\approx\!0.47$, yet the fine-tuned LLM RDM points in the
\emph{opposite direction} ($r\!=\!-0.88$) --- sign and contrast with the noise
ceiling are the load-bearing quantities, not the exact magnitude.
The 11D systems-level result (Section~\ref{sec:inversion}) provides stronger support
because its bootstrap CI for the fine-tuned comparison remains fully negative
($[-0.9899,\,-0.1178]$).

\textbf{(5b) RDM comparison in 11D systems-level space.}

\begin{center}
\small
\begin{tabular}{lcc}
\toprule
Comparison & 48D $r$ & 11D $r$ \\
\midrule
Brain vs.\ fine-tuned LLM  & $-0.6371$ & $-0.7539$ \\
Brain vs.\ pretrained LLM  & $-0.1023$ & $-0.9918$ \\
\bottomrule
\end{tabular}
\end{center}

As biological signal is denoised from 48 ROIs to 11 systems-level features, the
divergence \emph{amplifies} rather than shrinks.
The pretrained LLM at 11D ($r\!=\!-0.9918$) is the strongest cross-system result in
the study; this specific magnitude carries the uncertainty of a 3-point Pearson
correlation (1 df) and should be treated as indicative.
The robust directional claim is: as biological noise is removed, the negative
correlation between brain and LLM representational geometry becomes more extreme,
not less --- inconsistent with the hypothesis that the paradox is noise-driven.

\textbf{(5c) Valence-arousal axis alignment.}
We projected emotion centroids into the leading 2D PCA subspace and correlated each
axis with external V-A reference coordinates (Table~\ref{tab:va-coords}).

\vspace{0.4em}
\noindent\refstepcounter{table}\label{tab:va-coords}%
\begin{center}
\small
\begin{tabular}{lccc}
\toprule
Emotion & Valence & Arousal & Quadrant \\
\midrule
Afraid / Fear       & $-0.80$ & $0.85$ & negative, high arousal \\
Calm                & $+0.70$ & $0.15$ & positive, low arousal \\
Delighted / Happin. & $+0.85$ & $0.70$ & positive, moderate--high arousal \\
Depressed / Sadness & $-0.85$ & $0.25$ & negative, low arousal \\
Excited             & $+0.75$ & $0.90$ & positive, high arousal \\
Love                & $+0.85$ & $0.55$ & positive, moderate arousal \\
\bottomrule
\end{tabular}
\par\vspace{0.3em}
{\small\textbf{Table~\thetable:} External V-A reference coordinates (normalised
circumplex scale: valence $-1$ to $+1$, arousal $0$ to $1$).
The quadrant structure --- not exact decimal values --- is the
scientifically load-bearing element.}
\end{center}
\vspace{0.3em}

\begin{figure}[h!]
  \centering
  \begin{subfigure}[b]{0.48\textwidth}
    \includegraphics[width=\textwidth]{va_alignment_Brain-fMRI}
    \caption{Brain fMRI: poor VA alignment (Disparity$\!=\!0.80$, $R_\text{val}\!=\!0.31$).
             Embeddings (red crosses) deviate markedly from the V-A circumplex
             (blue circles).}
    \label{fig:va-brain}
  \end{subfigure}
  \hfill
  \begin{subfigure}[b]{0.48\textwidth}
    \includegraphics[width=\textwidth]{va_alignment_MPNet-Balanced}
    \caption{MPNet-Balanced (fine-tuned): excellent valence alignment
             (Disparity$\!=\!0.19$, $R_\text{val}\!=\!0.97$).
             Embeddings closely track the horizontal valence axis.}
    \label{fig:va-llm}
  \end{subfigure}

  \caption{\textbf{The Valence-Arousal Inversion (Finding 5c).}
           LLMs organise emotion primarily by valence; the brain primarily by arousal.
           The global representational priorities are swapped.}
  \label{fig:va-inversion}
\end{figure}

\vspace{0.4em}
\noindent\refstepcounter{table}\label{tab:va-alignment}%
\begin{center}
\small
\begin{tabular}{lcccc}
\toprule
System & PC1 axis & PC1 $r$ & PC2 axis & Permutation $p$ \\
\midrule
Fine-tuned LLM  & Valence & 0.98 & Arousal ($r\!=\!0.82$) & 0.013 \\
Pretrained LLM  & Valence & 0.97 & Arousal ($r\!=\!0.73$) & 0.009 \\
Human Brain     & Arousal & 0.96 & Valence ($r\!=\!0.31$) & 0.033 \\
\bottomrule
\end{tabular}
\par\vspace{0.3em}
{\small\textbf{Table~\thetable:} Global representational axis alignment
(Figure~\ref{fig:va-inversion}).
The axis priorities are swapped between LLMs and the brain.}
\end{center}
\vspace{0.3em}

Both LLM types are valence-dominant on PC1 ($r\!=\!0.97$--$0.98$);
the brain is arousal-dominant ($r\!=\!0.96$), with only weak valence alignment on PC2
($r\!=\!0.31$).
Because $K$ is small (5 brain emotions, 6 LLM emotions) and V-A coordinates are
category-level approximations, these results are presented as global-structure evidence
rather than primary proof.
The defensible claim is: \emph{the dominant global axis of the tested neural emotion
geometry aligns more strongly with an externally defined arousal reference than with
a valence reference, while the main ambiguity mechanism is supported independently
by centroid margin and classifier uncertainty.}

% ===========================================================================
\section{Discussion}
% ===========================================================================

\textbf{The shared geometric principle.}
Both LLM embedding spaces and the tested brain fMRI dataset exhibit the same general
relationship: geometric positioning relative to class prototypes predicts
representational certainty.
In LLMs this operates via centroid margin predicting logit margin; in the brain,
via centroid margin predicting classifier uncertainty.
These are complementary implementations of a geometric competition principle.
The mechanistic interpretation likely lies in the mathematics of distributed
high-dimensional representations: any system encoding discrete categories via
distributed activation will produce centroids surrounded by distributed instances,
generating a void, a belt, and an ambiguity gradient.

\textbf{Statistical hierarchy of evidence.}
The findings differ in evidentiary strength.
The \emph{strongest} evidence is the margin-uncertainty coupling: brain centroid margin
predicts classifier uncertainty (AUC$\,{=}\,0.81$ vs.\ 0.56 for raw distance),
collapsing under feature shuffling; LLM centroid margin predicts logit margin
($r\!=\!0.99$, permutation $p\!=\!0.00$).
\emph{Strong} evidence: the valence--arousal axis inversion, supported by independent
permutation tests on PC-axis correlations ($p\,{<}\,0.05$ for all three systems).
\emph{Supportive} evidence: the 11D RDM inversion, whose bootstrap CI for
fine-tuned LLMs is fully negative.
\emph{Descriptive evidence only}: the 48D RDM magnitude (3-point Pearson, wide CI)
and the void--belt density similarity (Wasserstein permutation $p\!=\!0.46$).
The primary evidence for global inversion is therefore the valence--arousal axis
analysis, not the 3-emotion RSA magnitude.

\textbf{The relational paradox: different priorities, not failure.}
The near-opposite RDM directionality does not indicate that either system is wrong.
LLMs trained on text corpora encode emotion co-occurrence statistics from natural
language --- ``fear'' and ``sadness'' share negative hedonic contexts and appear in
similar syntactic environments, placing their centroids near each other along the
primary valence axis.
The brain evolved to respond to environmental threats and rewards: fear involves
a physiological arousal response fundamentally different from the low-arousal
withdrawal of sadness, and separating states by activation intensity is most predictive
of appropriate behavioural response.
Neither system is incorrect; they represent emotion faithfully within their
respective operational domains.

\textbf{The low-rank emotion subspace.}
The 2.6\%-capacity finding has direct implications for efficient emotion representation.
Accuracy in 20D equals or exceeds 768D in all four tested variants.
For real-time affective computing or edge deployment, projecting to the
emotion-discriminative subspace is sufficient; larger representations may be
counterproductive without fine-tuning.

\textbf{Limitations.}
Our evidence covers one text emotion dataset (6 classes) and one brain fMRI dataset
(5 emotions, 40 subjects, one ROI atlas).
The 3-emotion cross-system label mapping (fear$\approx$afraid, happiness$\approx$delighted,
sadness$\approx$depressed) relies on conceptual equivalence and introduces non-trivial
uncertainty into all RSA comparisons.
The V-A analysis relies on externally assigned category-level coordinates rather than
participant-specific ratings; the global-axis findings should be interpreted as
supportive evidence, not as the main proof of the margin-based ambiguity mechanism.
All findings are observational; no causal evidence is established.

% ===========================================================================
\section{Conclusion}
% ===========================================================================

We have presented five findings characterising the geometry of emotion representations
in LLMs and human fMRI.
The core result is a paradox: a geometric competition principle for representational
ambiguity is observable across both systems --- LLMs via centroid margin predicting
logit margin; the brain via centroid margin predicting classifier uncertainty ---
yet global representational organisation is inverted: LLMs sort emotions by valence,
the brain by arousal.

Three findings have immediate practical value.
The low-rank emotion subspace (2.6\% of embedding capacity, 20 of 768 dimensions)
shows that efficient emotion encoding does not require large representations.
The margin-vs-distance contrast (AUC 0.81 vs.\ 0.56) provides a reusable diagnostic
distinguishing relational from radial geometric structure.
The V-A inversion suggests that training LLMs to distinguish arousal levels within the
same valence category may be necessary for biological alignment --- a hypothesis for
future experimental work.

More broadly, the coexistence of shared local geometry with globally inverted structure
suggests that cross-system alignment in affective representation is not a binary
question: two systems can share the same local competitive logic while differing
fundamentally in what they prioritise globally.

% ===========================================================================
\bibliographystyle{abbrvnat}
\bibliography{references}
% ===========================================================================

% ===========================================================================
% APPENDICES
% ===========================================================================
\appendix

\section{V-A Reference Coordinate Sensitivity Analysis}
\label{app:va-sensitivity}

The external V-A reference coordinates used in Section~\ref{sec:inversion} are
category-level approximations.
To demonstrate qualitative robustness:

\textbf{Step 1.} Baseline V-A table as in Table~\ref{tab:va-coords}.

\textbf{Step 2.} Add Gaussian noise to all valence and arousal coordinates at three
levels: SD\,=\,0.25, 0.50, and 0.75 (on the 1--9 scale).

\textbf{Step 3.} Repeat the PC-axis correlation analysis 1,000--10,000 times per noise
level.

\textbf{Step 4.} Report the proportion of perturbations for which the qualitative
result holds: PC1 of the brain remains arousal-dominant; PC1 of both LLM types remains
valence-dominant.

\textbf{Step 5.} Repeat using Spearman rank correlations.

The expected qualitative result is that arousal-dominant brain organisation and
valence-dominant LLM organisation are stable across perturbations reflecting the
quadrant structure of the circumplex model~\citep{posner2005}.
Code available at [anonymised repository].

% ---------------------------------------------------------------------------
\section{Phase 5 Full Results}
\label{app:phase5}
% ---------------------------------------------------------------------------

Phase 5 evaluates only the two fine-tuned final-layer models against the 4,038-sample
validation set.

\begin{center}
\small
\begin{tabular}{lcc}
\toprule
Metric & BGE-FT-Final & MPNet-FT-Final \\
\midrule
Total samples               & 4,038   & 4,038 \\
Overlap samples             & $\approx$76 (1.88\%) & $\approx$71 (1.76\%) \\
Dist-logit Pearson $r$      & 0.9572  & 0.9884 \\
Logit agreement in overlap  & 93.4\%  & 100\%  \\
Dominant overlap pair       & Happin./Love & Happin./Love \\
\bottomrule
\end{tabular}
\end{center}

The 100\% logit agreement for MPNet-FT-Final means every sample geometrically closer to
a competing centroid also has a higher raw logit score for that competing class.
Embedding geometry is the direct mechanistic basis for classifier confidence in this
model.

% ---------------------------------------------------------------------------
\section{Overlap Metrics: All Final-Layer Variants}
\label{app:overlap}
% ---------------------------------------------------------------------------

\begin{center}
\small
\begin{tabular}{lcccc}
\toprule
Variant & Overlap (\%) & Density Peak & Certainty Buffer & Void Onset \\
\midrule
MPNet-FT-Final   & 1.76  & 0.562 & $+1.875$ & $\approx$0.4375 \\
BGE-FT-Final     & 1.88  & 0.688 & $+1.750$ & $\approx$0.4375 \\
MPNet-Base-Final & 18.45 & 0.938 & $\phantom{+}0.000$ & $\approx$0.8125 \\
BGE-Base-Final   & 19.22 & 0.938 & $+0.125$ & $\approx$0.6875 \\
BGE-Base-Mid     & 28.14 &  ---  &    ---   &    ---    \\
MPNet-Base-Mid   & 16.92 &  ---  &    ---   &    ---    \\
\bottomrule
\end{tabular}
\end{center}

% ---------------------------------------------------------------------------
\section{fMRI Sample Size Justification}
\label{app:sample-size}
% ---------------------------------------------------------------------------

\begin{center}
\small
\begin{tabular}{lcc}
\toprule
Dataset & $N$ subjects & Venue \\
\midrule
Individual Brain Charting (IBC)         & 12 & Nature Scientific Data \\
Precision Functional Mapping (Children) & 12 & Dosenbach Lab / PMC \\
Natural Scenes Dataset (NSD)            & 8  & NSD Website \\
\textbf{This study (DS005700)}          & \textbf{40} & OpenNeuro \\
\bottomrule
\end{tabular}
\end{center}

With $N\!=\!40$, LOSO cross-validation leaves 39 training subjects per fold ---
well within the range where logistic regression on 48 features is well-determined.
The 5-class chance level is 0.20; our LOSO accuracy of 0.56
(95\% CI: 0.49--0.63, permutation $p<0.001$) represents an absolute lift of 0.36
above chance.

% ---------------------------------------------------------------------------
\section{48D$\to$11D Transformation: ROI-to-Network Mapping}
\label{app:11d}
% ---------------------------------------------------------------------------

\textbf{Anatomical lobe features (5):} Mean activation within
(1) frontal cortex ROIs,
(2) temporal cortex ROIs,
(3) parietal cortex ROIs,
(4) occipital cortex ROIs,
(5) cingulate/limbic ROIs.

\textbf{Functional network features (5):} Mean activation within
(6) Default Mode Network ROIs~\citep{li2014dmn} --- bilateral medial PFC, posterior
cingulate, angular gyrus;
(7) Salience Network ROIs~\citep{seeley2019} --- bilateral anterior insula, anterior
cingulate cortex;
(8) Central Executive Network ROIs~\citep{miller2001} --- bilateral dorsolateral PFC,
lateral parietal cortex;
(9) visual processing network ROIs --- bilateral occipital and fusiform regions;
(10) somatomotor network ROIs --- bilateral pre/postcentral gyrus.

\textbf{Neighbour context feature (1):} Mean absolute difference between adjacent ROI
activations within each anatomical region, capturing local inter-regional coordination.

% ---------------------------------------------------------------------------
\section{Image Experiments: Negative Result}
\label{app:image}
% ---------------------------------------------------------------------------

An extension to 120,000 affective image embeddings (CLIP-based) was explored.
Image embeddings did not yield clear cluster separation: silhouette scores remained
near zero and the void-belt density structure was absent.
The failure mode was insufficient cluster separation across emotion categories.
This negative result informs the scope of the geometric competition principle:
it is observable in text-based LLM embeddings and brain fMRI representations under the
systems and datasets tested, but was not apparent in CLIP image embeddings.
The modality (text vs.\ image), CLIP's training objective, and the emotion taxonomy
used are candidate explanations.

% ---------------------------------------------------------------------------
\section{Compute and Reproducibility}
\label{app:compute}
% ---------------------------------------------------------------------------

\textbf{Hardware.}
LLM embedding extraction: HPC cluster (GPU nodes).
Brain analysis: standard compute (CPU-bound statistical analyses).

\textbf{Code and data.}
Available at [anonymised repository].
The repository contains:
(1) embedding extraction scripts for all 8 LLM variants,
(2) all six brain analysis scripts,
(3) the 48D$\to$11D transformation code,
(4) the cross-system global behaviour comparison script,
(5) the V-A alignment code,
(6) \texttt{valence\_arousal\_reference\_table.csv},
(7) requirements and environment specification.

\textbf{Ethics statement.}
This study uses only publicly available, de-identified neuroimaging data from the
Neural MO dataset, OpenNeuro accession DS005700.
No new participants were recruited.
Informed consent and ethical approval were obtained by the original dataset creators;
this re-analysis does not require additional ethical review.

\end{document}
